\documentclass[12pt,a4paper]{article}
\usepackage[utf8]{inputenc}
\usepackage[T1]{fontenc}
\usepackage[spanish]{babel}
\usepackage[margin=2.5cm]{geometry}
\usepackage[hidelinks]{hyperref}

\setlength{\parindent}{0pt}
\setlength{\parskip}{6pt}

\title{Aplicación móvil: Nivel I\\Tutorial de Android Studio desde cero}
\author{Taller de Innovación y Emprendimiento Tecnológico\\Universidad Adolfo Ibáñez}
\date{15 de octubre de 2026}

\begin{document}
\maketitle

\section{Objetivo}

Esta guía corresponde a la calculadora versión 1.1.1, con Kotlin para Android y HTML, CSS y JavaScript para la interfaz web. Este tutorial explica cómo instalar Android Studio en Windows, descargar los archivos del proyecto, configurar las herramientas de compilación, generar una calculadora para Android e instalarla en un teléfono.

En Android, el archivo que se instala manualmente termina en \texttt{.apk}. La palabra \emph{app} significa aplicación; en este procedimiento no se genera un archivo con extensión \texttt{.app}.

La calculadora utiliza HTML para su estructura, CSS para su presentación y JavaScript para las operaciones. El proyecto completo que acompaña a esta guía incluye un contenedor Android escrito en Kotlin. Kotlin gestiona la ventana, el WebView y el ciclo de vida de Android. JavaScript realiza los cálculos y responde a los botones; HTML define la estructura y CSS define la apariencia. No hay archivos de código Java de la app. El JDK sigue siendo una herramienta de compilación necesaria para Kotlin.

\section{Material necesario}

\begin{itemize}
    \item Un ordenador Windows de 64 bits compatible con Android Studio.
    \item Conexión a Internet para descargar las herramientas y sincronizar el proyecto por primera vez.
    \item El proyecto completo de la calculadora, además del APK si ya existe uno.
    \item Un teléfono Android y un medio para transferir archivos, por ejemplo un cable USB.
\end{itemize}

Para generar el APK no necesitas un emulador ni un teléfono conectado. Tampoco necesitas instalar Git si descargas el repositorio como ZIP.

\section{Descargar e instalar Android Studio}

Las instrucciones de instalación y los requisitos del equipo se encuentran en la documentación oficial.\footnote{\url{https://developer.android.com/studio/install}}

\begin{enumerate}
    \item Abre la página oficial de Android Studio:
    \url{https://developer.android.com/studio}
    \item Descarga la versión estable para Windows. Para una instalación sencilla, elige el instalador con extensión \texttt{.exe}.
    \item Ejecuta el instalador y sigue sus indicaciones. Puedes conservar la carpeta de instalación propuesta.
    \item Abre Android Studio al terminar. Si pregunta si deseas importar una configuración anterior y es tu primera instalación, elige no importarla.
    \item En el asistente inicial, selecciona la instalación \textbf{Standard} si está disponible. Si utilizas una configuración personalizada, incluye el Android SDK.
    \item Revisa las licencias y acéptalas si estás de acuerdo. Espera a que termine la descarga de los componentes y pulsa \textbf{Finish}.
\end{enumerate}

El SDK reúne las herramientas y las API necesarias para construir aplicaciones Android. No es necesario crear un proyecto nuevo en este momento.

\section{Descargar el repositorio de GitHub}

Repositorio del taller:

\url{https://github.com/artemioadaysolvers/Aplicacion-Movil-Nivel-I}

\begin{enumerate}
    \item Abre el enlace anterior en tu navegador.
    \item En la página del repositorio, pulsa \textbf{Code}.
    \item Selecciona \textbf{Download ZIP}. Este procedimiento descarga una copia de los archivos de la rama seleccionada.\footnote{\url{https://docs.github.com/en/repositories/working-with-files/using-files/downloading-source-code-archives}}
    \item Busca el ZIP en las descargas del ordenador.
    \item Haz clic derecho y selecciona \textbf{Extraer todo}. Guarda los archivos en una carpeta local fácil de identificar.
\end{enumerate}

\subsection{Comprobar si están los archivos del proyecto}

Al revisar la rama \texttt{main} el 15 de octubre de 2026, el repositorio contenía \texttt{README.md} y \texttt{app-debug.apk}. El APK permite instalar la aplicación, pero no contiene el proyecto editable para volver a compilarlo con este procedimiento.

Para continuar con la compilación necesitas el código fuente completo. Puedes utilizar el archivo \texttt{CalculadoraSencillaKotlin\_AndroidStudio\_v1.1.1.zip}, entregado junto a este tutorial, y extraerlo en otra carpeta. Si posteriormente se incorpora el proyecto completo al repositorio, podrás obtenerlo directamente con \textbf{Download ZIP}.

La carpeta correcta del proyecto de esta calculadora contiene, entre otros, estos elementos:

\begin{verbatim}
CalculadoraSencillaKotlin/
    settings.gradle.kts
    build.gradle.kts
    gradle.properties
    gradlew
    gradlew.bat
    gradle/wrapper/gradle-wrapper.jar
    gradle/wrapper/gradle-wrapper.properties
    app/build.gradle.kts
    app/src/main/AndroidManifest.xml
    app/src/main/assets/www/
\end{verbatim}

Si solamente encuentras un APK y un README, todavía falta el proyecto completo. No utilices la opción de analizar un APK como sustituto de abrir el código fuente.

\section{Abrir el proyecto en Android Studio}

\begin{enumerate}
    \item En la pantalla inicial de Android Studio, pulsa \textbf{Open}. Si ya tienes otro proyecto abierto, utiliza \textbf{File > Open}.
    \item Busca la carpeta extraída del proyecto completo.
    \item Selecciona \textbf{CalculadoraSencillaKotlin}, la carpeta que contiene \texttt{settings.gradle.kts}, \texttt{build.gradle.kts} y \texttt{gradlew.bat}.
    \item Abre el proyecto y espera a que Android Studio termine de reconocer sus archivos.
\end{enumerate}

No selecciones únicamente la carpeta \texttt{app} ni la carpeta \texttt{www}. Tampoco debes crear un proyecto vacío y copiar los archivos dentro de él.

\section{Configurar el SDK y Java}

Los pasos siguientes corresponden al proyecto completo versión 1.1.1. Sus versiones configuradas son Android Gradle Plugin 8.13.2, Kotlin 2.3.20, Gradle 8.13 y Android SDK 36. El complemento de Android 8.13 requiere como mínimo JDK 17.\footnote{\url{https://developer.android.com/build/releases/agp-8-13-0-release-notes}}

\subsection{Instalar los componentes del SDK}

\begin{enumerate}
    \item Abre \textbf{Tools > SDK Manager}.
    \item En \textbf{SDK Platforms}, selecciona \textbf{Android SDK Platform 36}, correspondiente a Android 16, si todavía no está instalada.
    \item En \textbf{SDK Tools}, comprueba que estén instaladas las \textbf{Android SDK Platform-Tools}.
    \item Instala también \textbf{Android SDK Build-Tools 35.0.0} si Android Studio lo solicita para este proyecto. Si necesitas seleccionar una versión específica, activa la vista de detalles de los paquetes.
    \item Pulsa \textbf{Apply}, revisa las licencias y espera a que termine la instalación.
\end{enumerate}

La plataforma utilizada para compilar no tiene que coincidir con la versión del teléfono. Este proyecto establece Android 6.0 como versión mínima; puede instalarse en un teléfono con Android 15 aunque se compile con el SDK 36.\footnote{\url{https://developer.android.com/studio/intro/update#sdk-manager}}

\subsection{Seleccionar el JDK para Gradle}

\begin{enumerate}
    \item Abre \textbf{File > Settings}.
    \item Entra en \textbf{Build, Execution, Deployment > Build Tools > Gradle}.
    \item En \textbf{Gradle JDK}, selecciona \textbf{JDK 17}. Si no está disponible, utiliza \textbf{Download JDK}, elige la versión 17 y completa la descarga.
    \item Pulsa \textbf{Apply} y después \textbf{OK}.
\end{enumerate}

Esta configuración elige Java para la compilación del proyecto; no exige cambiar el Java con el que se inicia Android Studio.\footnote{\url{https://developer.android.com/build/jdks#gradle-jdk}}

\section{Sincronizar el proyecto}

\begin{enumerate}
    \item Si aparece \textbf{Sync Now}, pulsa esa opción. También puedes utilizar \textbf{File > Sync Project with Gradle Files}.
    \item Mantén Internet conectado mientras Gradle descarga los componentes necesarios.
    \item Si aparece \textbf{Install missing SDK packages}, revisa e instala los componentes indicados.
    \item Espera a que termine la sincronización. Si falla, revisa el mensaje de error antes de continuar.
\end{enumerate}

El proyecto incluye Gradle Wrapper, que selecciona y descarga la versión de Gradle indicada en sus archivos. No hace falta instalar Gradle por separado. Para esta calculadora tampoco se requieren Node.js, npm ni Capacitor.

Conserva las versiones configuradas del proyecto durante esta primera compilación. Una sugerencia para actualizar herramientas no significa necesariamente que debas cambiarlas para generar el APK.

\section{Generar el APK de la calculadora}

\begin{enumerate}
    \item Abre \textbf{View > Tool Windows > Build Variants} y selecciona \textbf{debug} para el módulo \textbf{app}.
    \item En el menú superior, pulsa \textbf{Build}.
    \item Selecciona \textbf{Generate Bundle(s) / APK(s)} y después \textbf{Generate APK(s)}. En otras versiones del IDE puede aparecer como \textbf{Build Bundle(s) / APK(s) > Build APK(s)}.
    \item Espera a que finalice el proceso. La primera compilación puede tardar más porque descarga herramientas y prepara archivos.
    \item Comprueba que el resultado indique \texttt{BUILD SUCCESSFUL}.
\end{enumerate}

Elige APK para instalar manualmente en el teléfono. Un Android App Bundle, con extensión \texttt{.aab}, corresponde a otro procedimiento de distribución. Para esta práctica, el APK de la variante \texttt{debug} ya queda firmado con una clave de desarrollo; no hace falta usar el asistente de firma de distribución.\footnote{\url{https://developer.android.com/build/build-for-release}}

\section{Encontrar el archivo correcto}

Al terminar, pulsa \textbf{Locate} en la notificación de Android Studio si aparece. También puedes buscar el resultado desde el Explorador de archivos de Windows, dentro de la carpeta del proyecto:

\begin{verbatim}
app/build/outputs/apk/debug/app-debug.apk
\end{verbatim}

El archivo que debes copiar al teléfono se llama \textbf{app-debug.apk}. Si aparece otro archivo llamado \texttt{app-debug-androidTest.apk}, ese es un paquete de pruebas y no es la calculadora que debes abrir.

\texttt{BUILD SUCCESSFUL} confirma que la compilación terminó correctamente. Después debes instalar y abrir la aplicación para comprobar su funcionamiento.

\section{Instalar y abrir la aplicación en el teléfono}

\begin{enumerate}
    \item Conecta el teléfono al ordenador mediante USB y elige \textbf{Transferencia de archivos} si el teléfono lo solicita.
    \item Copia \texttt{app-debug.apk} a la carpeta \textbf{Descargas} del teléfono.
    \item En el teléfono, abre \textbf{Archivos} y entra en \textbf{Descargas}.
    \item Toca \texttt{app-debug.apk}. Si Android solicita permiso para instalar desde la aplicación Archivos, habilítalo para ese origen y vuelve al instalador.
    \item Pulsa \textbf{Instalar}. Si actualizas la misma calculadora y el APK utiliza la misma clave de firma, puede aparecer \textbf{Actualizar}.
    \item Espera a que termine y pulsa \textbf{Abrir}.
\end{enumerate}

Los nombres de los menús pueden variar entre fabricantes. En esta calculadora, el nombre visible es \textbf{Calculadora} y el icono tiene una calculadora blanca sobre fondo naranja.

Para actualizarla, recompila el proyecto corregido y copia el APK recién generado. Si la versión anterior y la nueva utilizan claves de firma distintas, Android no permitirá sustituir una por otra como actualización directa.

\section{Comprobar el funcionamiento}

Después de abrir la calculadora, realiza estas pruebas:

\begin{itemize}
    \item Suma: $2+3=5$.
    \item Resta: $9-4=5$.
    \item Multiplicación: $6\times7=42$.
    \item División: $8\div2=4$.
    \item Decimales: $0{,}1+0{,}2=0{,}3$.
    \item División entre cero: debe aparecer un mensaje de error sin cerrar la aplicación.
\end{itemize}

El botón \textbf{AC} borra la operación. La calculadora funciona sin conexión una vez instalada.

\section{Problemas frecuentes}

\subsection{La aplicación no aparece en el teléfono}

Comprueba que instalaste \texttt{app-debug.apk}, sin \texttt{androidTest} en el nombre. Busca también \textbf{Calculadora} en la lista de aplicaciones del teléfono. Si todavía tienes abierta la confirmación de instalación, utiliza su botón \textbf{Abrir}.

\subsection{No se encuentra el SDK o hay un error de Java}

Si aparece \texttt{SDK location not found}, revisa la ubicación del Android SDK en Android Studio. Si falta la plataforma 36, instálala mediante \textbf{SDK Manager}. Si el error menciona la versión de Java, selecciona JDK 17 en los ajustes de Gradle y sincroniza de nuevo.

\subsection{Aparece una advertencia sobre funciones obsoletas}

Una advertencia de funciones obsoletas de Gradle no equivale por sí sola a una compilación fallida. Comprueba el resultado final: \texttt{BUILD SUCCESSFUL} indica éxito; \texttt{BUILD FAILED} requiere revisar el error correspondiente.

\subsection{La calculadora se cierra al abrirse}

Comprueba primero que compilaste el proyecto corregido versión 1.1.1 y que instalaste su APK nuevo. Si continúa cerrándose, obtén el registro del teléfono:

\begin{enumerate}
    \item Activa las \textbf{Opciones de desarrollador} y la \textbf{Depuración USB} del teléfono. La ubicación depende del fabricante; normalmente se habilitan pulsando varias veces el número de compilación en la información del dispositivo.
    \item Conecta el teléfono mediante USB y autoriza la depuración para tu ordenador.\footnote{\url{https://developer.android.com/studio/run/device}}
    \item En Android Studio, abre \textbf{View > Tool Windows > Logcat} y selecciona el teléfono conectado.
    \item En la búsqueda de Logcat puedes escribir \texttt{package:mine is:crash} para filtrar cierres de las aplicaciones del proyecto abierto.
    \item Abre la calculadora otra vez. Copia el error y su traza para identificar el punto que falla.\footnote{\url{https://developer.android.com/studio/debug/logcat}}
\end{enumerate}

\section{Modificar la calculadora y volver a compilar}

En el proyecto completo, los archivos web están en:

\begin{verbatim}
app/src/main/assets/www/
\end{verbatim}

\begin{itemize}
    \item \texttt{index.html}: estructura de la pantalla y botones.
    \item \texttt{styles.css}: presentación y tamaños.
    \item \texttt{engine.js}: operaciones de la calculadora.
    \item \texttt{app.js}: interacción de los botones y el teclado.
\end{itemize}

Guarda los cambios, genera otra vez el APK y transfiere el archivo nuevo al teléfono. No basta con modificar los archivos del ordenador para que la aplicación instalada cambie.

\section{Usar este archivo en Overleaf}

Crea un proyecto mediante \textbf{New Project > Blank Project} y pega este código en \texttt{main.tex}. También puedes subir el archivo \texttt{.tex} al proyecto. No se necesitan imágenes, archivos externos ni una plantilla adicional.\footnote{\url{https://docs.overleaf.com/getting-started/your-first-project}}

Este material se entrega como código fuente LaTeX, sin generar un PDF.

\end{document}
